\chapter{Introducción}
\label{cap:introduccion}

\section{Contexto, antecedentes y motivación}
En el ámbito de la investigación biomédica, el análisis de imágenes de resonancia magnética, o Imagen por Resonancia Magnética (IRM), constituye una herramienta fundamental para el estudio de estructuras anatómicas y el análisis de su evolución. En particular, en estudios preclínicos sobre modelos animales, como el cerebro de ratón, la resonancia magnética permite observar la evolución de distintas patologías, proporcionando información relevante para la investigación en neurociencia y el análisis de tratamientos. \citep{bioimagenManualRMImageJ2024,saito2024preclinicalMRI}

El procesamiento de este tipo de imágenes, conocido como \textbf{segmentación}, consiste en delimitar la región anatómica o patológica de interés, denominada \textbf{región de interés}, con el objetivo de cuantificar las características del órgano o de la lesión, así como estimar su volumen. Estas tareas resultan esenciales para garantizar la reproducibilidad de los experimentos y la comparabilidad de los resultados entre distintos estudios y sujetos, especialmente en aquellos casos en los que se analiza la evolución temporal de una patología o se está evaluando un posible tratamiento. \citep{bioimagenManualRMImageJ2024}

Tradicionalmente, estas tareas se han realizado de forma manual: el flujo de trabajo habitual consiste en analizar cada uno de los cortes de la resonancia (aproximadamente entre 20 y 50 cortes por sujeto), identificar las regiones de interés y delimitarlas manualmente mediante herramientas de selección específicas, a partir de las cuales se procede posteriormente al cálculo de magnitudes de interés.

Aunque este procedimiento permite obtener resultados experimentales válidos, presenta diversas limitaciones.
Se trata de un proceso laborioso que consume una gran cantidad de tiempo, especialmente en estudios con un elevado número de sujetos. Esta carga de trabajo supone una inversión significativa de tiempo por parte de los investigadores, lo que puede repercutir en los tiempos de desarrollo de la investigación.
Por otro lado, la delimitación manual de ciertas regiones, como las lesiones isquémicas, presenta un grado de ambigüedad, ya que sus fronteras no siempre están claramente definidas. Esto puede introducir variabilidad entre anotaciones en función del criterio del profesional. \citep{bioimagenManualRMImageJ2024,veigaCanuto2022interobserverMRI,liu2022interobserverSegmentation}

A estas limitaciones se suma la necesidad de analizar conjuntos de imágenes cada vez más extensos y de obtener medidas cuantitativas de forma consistente, así como la demanda de análisis detallados y precisos, lo que dificulta la escalabilidad de los procesos manuales.


En este contexto, surge la necesidad de desarrollar herramientas que permitan agilizar y automatizar de manera supervisada estas tareas, mejorando la eficiencia del proceso y aumentando la consistencia y reproducibilidad de los resultados.

El avance reciente de las técnicas de inteligencia artificial aplicadas a imagen ha demostrado su capacidad para identificar patrones complejos en datos visuales, abriendo la puerta a su aplicación en entornos biomédicos. No obstante, la integración de este tipo de soluciones en contextos reales de trabajo requiere no solo obtener resultados precisos, sino también garantizar que dichos resultados sean interpretables y coherentes desde el punto de vista del dominio de aplicación. \citep{ronneberger2015unet}

Este Trabajo de Fin de Grado, en colaboración con el centro ICTS BioImagen Complutense, pretende dar respuesta a una necesidad real observada en el entorno de investigación. A partir del conocimiento adquirido, el análisis del estado del arte y la revisión de la bibliografía, se plantea el desarrollo de una herramienta basada en modelos supervisados que permita contribuir a la mejora del proceso de análisis de imágenes de resonancia magnética en investigación preclínica, proponiendo una solución que reduzca la carga de trabajo manual y mejore la consistencia de los resultados.
La motivación de este trabajo no se limita únicamente a la implementación de un modelo automático, sino que se centra en el desarrollo de una herramienta que pueda integrarse dentro del flujo de trabajo existente, aportando valor como sistema de apoyo a la decisión, combinando la automatización con la supervisión por parte del experto, manteniendo el control sobre los resultados finales.

\section{Objetivos}

El objetivo del Trabajo de Fin de Grado es desarrollar una herramienta informática de apoyo al análisis de imágenes biomédicas preclínicas obtenidas mediante IRM. El trabajo desarrollado se centra en dos tareas de segmentación: la delimitación automática del cerebro y la segmentación de regiones correspondientes a isquemia cerebral.

La herramienta no pretende sustituir el criterio del investigador experto, sino reducir la carga asociada a la segmentación manual y proporcionar resultados reproducibles que puedan ser revisados, interpretados y utilizados como apoyo dentro del flujo de trabajo habitual del centro ICTS BioImagen Complutense.


\subsection{Objetivo general}
Desarrollar y evaluar un prototipo basado en técnicas de visión por computador e inteligencia artificial capaz de segmentar automáticamente, en imágenes de resonancia magnética preclínica, el área correspondiente al cerebro y las regiones afectadas por isquemia cerebral, proporcionando además información cuantitativa asociada a dichas segmentaciones.

\subsection{Objetivos específicos}
Para alcanzar el objetivo general se han definido los siguientes objetivos específicos:

\begin{itemize}
    \item Analizar el contexto biomédico del problema y el flujo de trabajo seguido en ICTS BioImagen Complutense para el análisis de imágenes de IRM preclínica asociadas a estudios de isquemia cerebral.
    \item Estudiar el proceso manual de segmentación del cerebro y de las regiones isquémicas, identificando sus principales limitaciones en términos de tiempo, reproducibilidad y dependencia del criterio del investigador.
    \item Organizar y preparar un conjunto de datos formado por imágenes reales de resonancia magnética, regiones de interés y máscaras binarias correspondientes exclusivamente a cerebro e isquemia cerebral.
    \item Diseñar un proceso de preprocesamiento que permita adaptar las imágenes originales al entrenamiento de modelos supervisados, incluyendo tareas de conversión de formatos, normalización, filtrado de cortes no informativos y ajuste dimensional.
    \item Entrenar un primer modelo de segmentación basado en una arquitectura U-Net para delimitar automáticamente el área del cerebro en las imágenes de entrada.
    \item Entrenar un segundo modelo de segmentación basado en una arquitectura U-Net para identificar las regiones de isquemia cerebral presentes en las imágenes.
    \item Utilizar la segmentación del cerebro como información auxiliar para restringir el análisis de isquemia a la región anatómica relevante y reducir detecciones fuera del área cerebral.
    \item Implementar mecanismos de cuantificación que permitan calcular medidas derivadas de las segmentaciones obtenidas, especialmente el volumen del cerebro y el volumen de la región isquémica.
    \item Evaluar el rendimiento de ambos modelos mediante métricas cuantitativas de segmentación y análisis visual de las predicciones generadas.
    \item Valorar posibles líneas de ampliación futura, ofreciendo un primer desarrollo de un software que permite entrenar modelos para otros órganos y patologías.
\end{itemize}

\section{Plan de trabajo}
El desarrollo de este Trabajo de Fin de Grado se ha organizado de forma progresiva, partiendo de una fase inicial de comprensión del problema biomédico hasta la implementación y evaluación del sistema propuesto. Dado que el proyecto combina aspectos de ingeniería del software, procesamiento de imagen médica, aprendizaje automático y validación experimental, resultó necesario estructurar el trabajo en distintas fases claramente diferenciadas.

Además, al tratarse de un proyecto desarrollado en colaboración con el centro ICTS BioImagen Complutense, el plan de trabajo estuvo condicionado por las necesidades detectadas durante las reuniones de seguimiento y las visitas realizadas al centro. Estas sesiones permitieron orientar las decisiones técnicas hacia un caso de uso real, definir los requisitos funcionales de la herramienta y adaptar el desarrollo a las características del conjunto de datos disponible.

En los siguientes apartados se describe, en primer lugar, el proceso de seguimiento llevado a cabo mediante actas de reunión y, posteriormente, las principales fases de desarrollo seguidas hasta la obtención del prototipo final.

\subsection{Seguimiento del proyecto y reuniones con los tutores}
A lo largo del desarrollo del proyecto se mantuvo un proceso de seguimiento continuado mediante reuniones periódicas con los tutores, tanto de manera presencial como telemática. Con el objetivo de tener registro de la evolución del desarrollo y las decisiones tomadas, se realizaron actas de reunión después de cada encuentro.

Estas actas tuvieron un papel especialmente importante durante las primeras fases del proyecto, ya que permitieron documentar la definición inicial del problema, el contexto de la aplicación y las necesidades del usuario final. Además, en las actas se recogieron los aspectos teóricos y técnicos guiados por los profesores, tanto en el ámbito de la inteligencia artificial como en el de la aplicación biomédica, donde los alumnos también tuvieron una formación básica. 

Durante las visitas al centro ICTS BioImagen Complutense, las actas permitieron recoger información esencial sobre el flujo de trabajo seguido por los investigadores. Se analizaron ejemplos reales de imágenes y se estudiaron diversas aplicaciones de las IRM tratadas en el centro. Además, se adquirieron nociones básicas sobre el funcionamiento físico de los equipos de resonancia magnética, aprendiendo a diferenciar zonas de hiperintensidad e hipointensidad, claves para la elaboración posterior del dataset.

Las actas también reflejan la evolución técnica del proyecto. A medida que avanzó el desarrollo, se tomaron decisiones relacionadas con la estructura del conjunto de datos, la eliminación de cortes poco informativos, la preparación de los notebooks de entrenamiento y la mejora del modelo de segmentación de isquemia. En particular, se decidió utilizar la segmentación cerebral como paso previo para restringir el área de análisis del modelo de isquemia al tejido cerebral, con el objetivo de reducir falsos positivos y mejorar la coherencia anatómica de las predicciones.

Por último, el seguimiento de las actas permitió incorporar nuevas necesidades detectadas durante el desarrollo, como el uso de notebooks a través de Google Colab, facilitando la accesibilidad del software a los usuarios finales, que serán los investigadores del centro ICTS BioImagen Complutense.

\subsection{Fases de desarrollo}
El desarrollo del trabajo se ha organizado en varias fases. 
En primer lugar, se realizó una fase de toma de contacto con el problema biomédico mediante reuniones con los directores del trabajo y visitas al centro ICTS BioImagen Complutense. En esta fase se estudió el proceso de adquisición de IRM, la generación de regiones de interés y la forma en la que los investigadores obtienen actualmente medidas cuantitativas a partir de las segmentaciones manuales.

Posteriormente, se abordó la preparación del conjunto de datos. Se seleccionó como caso de estudio la segmentación de lesiones de isquemia cerebral, debido tanto a la disponibilidad de datos proporcionados por el ICTS BioImagen Complutense como a su relevancia dentro de las líneas de investigación actuales del centro. En esta etapa se llevó a cabo la adaptación de las imágenes al formato requerido por los modelos de aprendizaje automático.

A continuación, se diseñó e implementó el pipeline de segmentación automática, compuesto por dos modelos: uno orientado a la segmentación del cerebro y otro a la segmentación de la lesión isquémica. Finalmente, se realizó la evaluación de los resultados obtenidos mediante métricas de segmentación y cuantificación volumétrica, complementando este análisis con técnicas de explicabilidad que permiten estudiar el comportamiento interno de los modelos.

\section{Estructura de la memoria}
La memoria se organiza de la siguiente forma. En el Capítulo 2 se presenta el estado de la cuestión, incluyendo los conceptos biomédicos necesarios, la visión por computador aplicada a imagen médica, las arquitecturas de segmentación, el aprendizaje por transferencia y la inteligencia artificial explicable. En el Capítulo 3 se describe la metodología seguida, desde la construcción del dataset y el preprocesamiento de imágenes hasta la arquitectura de los modelos, la cuantificación volumétrica y la integración de explicabilidad. En el Capítulo 4 se presentan y discuten los resultados obtenidos para la segmentación cerebral y la segmentación de isquemia. Finalmente, el Capítulo 5 recoge las conclusiones del trabajo y las principales líneas de trabajo futuro.